% !TeX root = ../常微分方程笔记.tex
% 由主文档通过 \input 载入；请编译主文档。
\setcounter{chapter}{1} % 保留原书章号。
\chapter{线性代数预备知识}
\label{app:ode-algebra}

\section{解空间、基与线性映射}
函数空间中的线性无关与向量空间完全相同：若 $\sum_{j=1}^n c_jy_j(t)\equiv0$ 只允许所有常数 $c_j=0$，则函数族线性无关。系数必须为常数；允许系数依赖 $t$ 会改变问题。

线性算子 $L$ 的齐次解集是核 $\ker L$。若 $L[y_p]=g$，则 $L[y]=g$ 的解集为 $y_p+\ker L$。有限维情形中，线性映射一旦同时单射、满射，就给出两个空间的维数相同；这正是用初值映射证明线性微分方程解空间维数的依据。

\section{特征值、Jordan 标准形与矩阵指数}
$\lambda$ 为矩阵 $A$ 的特征值，是指存在非零向量 $v$ 使 $Av=\lambda v$。它在特征多项式中的重数为代数重数，而 $\ker(A-\lambda I)$ 的维数为几何重数。后者不超过前者；所有特征值均有足够独立特征向量时，矩阵可对角化。

在复数域上，任意方阵都相似于 Jordan 标准形。每块为 $J=\lambda I+N$，其中 $N$ 只有紧邻对角线的上方可能为一，且为幂零矩阵。长度为 $m$ 的广义特征向量链满足
\[
 (A-\lambda I)v_1=0,\qquad (A-\lambda I)v_{j+1}=v_j.
\]
相应解包含 $e^{\lambda t}$ 乘次数至多 $m-1$ 的多项式。半单特征值是指其全部 Jordan 块均为一阶，不要求该特征值只有一重。

\section{范数、二次型与正定矩阵}
欧氏范数对应的诱导矩阵范数定义为 $\|A\|=\sup_{\|x\|=1}\|Ax\|$，因而 $\|Ax\|\leq\|A\|\|x\|$、$\|AB\|\leq\|A\|\|B\|$。这两条不等式支持矩阵指数和误差估计。

实对称矩阵 $P$ 正定是指 $x^{\mathsf T}Px>0$ 对所有 $x\ne0$ 成立。由正交对角化，若最小、最大特征值为 $m,M$，则
\[
 0<m\|x\|^2\leq x^{\mathsf T}Px\leq M\|x\|^2\quad(x\ne0).
\]
因此正定二次型同时控制状态范数并具有紧子水平集。若沿系统还有 $\dot V\leq-c\|x\|^2$，则 $\dot V\leq-(c/M)V$，从而 $\|x(t)\|\leq\sqrt{M/m}\,e^{-ct/(2M)}\|x(0)\|$。

